\ifdefined\MOneSmall
\documentclass[acmsmall,screen]{acmart}
\else
\documentclass[manuscript,screen,review]{acmart}
\fi
\usepackage{booktabs}
\usepackage{enumitem}
\usepackage[htt]{hyphenat}
\usepackage{seqsplit}
\setcopyright{none}
\settopmatter{printacmref=false}
\acmJournal{TOMPECS}
\acmDOI{}
\acmYear{2026}
\title{LLM-Serving Scheduler Ranking Portability Across Workloads, Operating Regions, and Metrics}
\author{Soroush Vahidi}
\email{sv96@njit.edu}
\affiliation{\institution{New Jersey Institute of Technology}\city{Newark}\state{New Jersey}\country{USA}}
\begin{document}
\begin{abstract}
Scheduler rankings are conditional on the workload, operating region,
metric and execution model used to obtain them. We present a reproducible
comparison methodology and an empirical study of a frozen, trace-driven
LLM-serving simulator corpus: three released trace sources, 120 windows,
six operating regions and 13 executed policies, including an 11-policy
primary panel. The 9,360 configurations each have two deterministic
verification passes, not independent stochastic repetitions. For the
primary arrival-normalized weighted-goodput metric, cross-source Kendall
correlations range from 0.551 to 1.000. Of 990 pairwise-policy comparison
records, 36 satisfy the saved practical-margin and nominal interval criteria for
a reversal; they span five policy pairs and all involve BurstGPT.
A post-campaign cross-metric analysis yields a median correlation of 0.419
and deterministic top-choice agreement of 31.9\%, while a post-hoc comparison finds overlap of maximizing policy sets in
57.8\% of conditions. Independent numerical reproduction confirms the saved
intervals, but their coverage and the historical FDR claim are not established. These characterize the historical
simulator, not native serving rankings. Retrospective audits of the legacy
GPU comparison and subsequent CPU conformance work leave faithful physical
ranking transfer unverified. The results support reporting comparative
rankings with their workload, metric, uncertainty and implementation boundaries.
\end{abstract}
\ccsdesc[500]{Computing methodologies~Simulation evaluation}
\ccsdesc[300]{Computer systems organization~Cloud computing}
\keywords{LLM serving, scheduling, ranking portability, trace-driven simulation, performance evaluation}
\maketitle
\input{revised/sections/introduction}
\input{revised/sections/related_work}
\input{revised/sections/research_questions}
\input{revised/sections/model}
\input{revised/sections/workloads}
\input{revised/sections/within}
\input{revised/sections/portability}
\input{revised/sections/metrics}
\input{revised/sections/robustness}
\input{revised/sections/limitations}
\input{revised/sections/conclusion}
\section*{Artifact availability and declarations}\label{sec:artifact}
The review repository preserves the source, configurations, frozen outcome
matrix, canonical analyses and the M1 and M2 claim audits. The dataset package
\texttt{dataset/v1.1.0} includes per-execution outcomes, analysis JSON and
reproducible tabular views; it does not include a full dump of the legacy
physical request records. Historical Hugging Face and Zenodo release
references require a current accessibility and version check before submission.
The working repository is private; this draft makes no claim that its new
review files are publicly available. Upstream trace redistribution remains
subject to the original terms.

The author declares no competing interests. Historical acknowledgments
identify support from Google Cloud Research Credits, the CloudRift AI Builder
Grant and the Cohere Labs Catalyst Grant Program, and acknowledge Professor
Ioannis Koutis and NJIT's Wulver/ARCS resources. Funding and contribution
statements require author confirmation for the final submission. The author
is responsible for conceptualization, methodology, software, validation,
analysis, data curation, visualization and writing. AI-assisted preparation
was used historically and Codex assisted these M1/M2 revisions; the revised
claims and text remain pending detailed author review and approval.

\bibliographystyle{ACM-Reference-Format}
\bibliography{revised/references}
\end{document}
